\section{综合例题：同一近场模型从场积分走到探测计数}
\label{sec:feqo-nearfield-counts-comprehensive}

前面的章节分别建立了场、电子能量格、自由传播、开放演化与测量语言。
下面固定一套装置：两段经典单频近场之间插入自由漂移与相位噪声，末端用有限能量分辨率的探测器记录电子。
各步始终沿用同一个电子态，因而可以区分场的相位、电子相干性、读出脉冲与计数噪声各自的作用。
第一段和第二段的耦合记为 $\beta_1,\beta_2$；这里的 $\beta_a$ 是无量纲复耦合，不是逆温度。

\Needspace{8\baselineskip}
\begin{exbox}[第一步：从结构近场算出复耦合]
设结构产生的纵向单频近场在电子轴线上为
\[
 E_{\omega,z}(z)=E_a\exp\!\left(-\frac{z^2}{2L^2}\right)\ee^{\ii \kappa z},
 \qquad L>0.
\]
这是结构近场的轴线剖面，横向场和结构边界满足前文的 Maxwell 问题。电子以中心速度 $v_0$ 穿过近场，求其复耦合及相位失配对耦合强度的影响。

\textbf{解。} 采用前文的电荷和场幅约定，将剖面代入纵向场积分。高斯 Fourier 积分给出
\begin{align*}
 \beta&=\frac{q_e}{\hbar\omega}\int_{\mathbb R}
 E_{\omega,z}(z)\ee^{-\ii \omega z/v_0}\,\dd z\\
 &=\frac{q_e E_a\sqrt{2\pi}L}{\hbar\omega}
 \exp\!\left[-\frac{L^2}{2}\left(\kappa-\frac{\omega}{v_0}\right)^2\right].
\end{align*}
故 $\beta=g \ee^{\ii \phi}$ 中的 $g$ 随相位失配呈高斯衰减，
\[
 g=\frac{|q_e E_a|\sqrt{2\pi}L}{\hbar\omega}
 \ee^{-L^2(\kappa-\omega/v_0)^2/2},
 \qquad \phi=\arg(q_eE_a).
\]
负电子电荷带来的相位也包含在 $\phi$ 中。后文估计 $g$ 时，只有在
$L,\kappa,v_0$ 已知的条件下，才能反解场幅 $|E_a|$。
\end{exbox}

\begin{exbox}[第二步：检查能量格近似，再计算第一次散射]
确定各能量通道的动量，检查何时可用同一个复耦合计算第一次散射的振幅与能谱。

\textbf{解。} 结构可以交换动量，有限空间包络本身不指定唯一动量转移。在单频且频宽可忽略的散射模型中，先固定能量格，再由色散关系反解动量：
\[
 E_\ell=E_0+\ell\hbar\omega,\qquad
 p_\ell=\frac1c\sqrt{E_\ell^2-m^2c^4},\qquad
 k_\ell=\frac{p_{\ell+1}-p_\ell}{\hbar}.
\]
在实际占据的有限能级窗内要求 $E_\ell>mc^2$，并要求各通道的失配变化
在近场长度上不能积累显著相位，例如
\[
 \max_{|\ell|\le M}L|k_\ell-\omega/v_0|\ll1.
\]
同时须控制速度变化、场频宽及不同能级的耦合变化，使第一步的同一个 $\beta$ 可用于所有相关边。
仅有 $\hbar\omega/E_0\ll1$ 并不足以保证这些条件。

在已经满足这些条件的理想双向能量格 $\ell^2(\mathbb Z)$ 中，定义
\[
 b|\ell\rangle=|\ell-1\rangle,\qquad
 b^\dagger|\ell\rangle=|\ell+1\rangle,\qquad
 bb^\dagger=b^\dagger b=I.
\]
因此此处 $b$ 不是具有真空底态的玻色湮灭算符。近场散射为
\[
 S(\beta)=\exp(\beta b^\dagger-\beta^*b),\qquad
 S(\beta)|0\rangle=\sum_{\ell\in\mathbb Z}a_\ell|\ell\rangle,
 \qquad a_\ell=J_\ell(2g)\ee^{\ii \ell\phi}.
\]
这一步同时给出振幅与能谱：
\[
 P_\ell^{(1)}=J_\ell^2(2g),\qquad
 \sum_\ell P_\ell^{(1)}=1,
 \qquad P_0^{(1)}=1-2g^2+O(g^4),\quad
 P_{\pm1}^{(1)}=g^2+O(g^4).
\]
由 $J_{-\ell}^2=J_\ell^2$ 得平均净能级变化为零。
在相位表示中，波函数为
\[
 \psi(\theta)=\frac{1}{\sqrt{2\pi}}
 \ee^{\beta \ee^{-\ii \theta}-\beta^*\ee^{\ii \theta}},\qquad
 N=\ii\partial_\theta.
\]
Parseval 恒等式于是给出
\[
 \langle N^2\rangle
 =\int_0^{2\pi}|\partial_\theta\psi|^2\,\dd\theta
 =\frac{1}{2\pi}\int_0^{2\pi}4g^2\cos^2(\phi-\theta)\,\dd\theta
 =2g^2.
\]
电子能谱展宽了，但单次能量测量仍完全看不到 $\phi$。
\end{exbox}

\begin{exbox}[第三步：自由漂移怎样保留能谱并改变相干性]
电子穿过两近场之间的固定距离 $Z_d$。求漂移后的密度矩阵，并比较其对角元与非对角元。

\textbf{解。} 固定距离的传播相位由同一能量格的 $p(E)$ 决定。
在 $E_0$ 附近展开动量：
\[
 p_\ell=p_0+\frac{\ell\hbar\omega}{v_0}
 -\frac{(\ell\hbar\omega)^2}{2m_\parallel v_0^3}+R_\ell^{(p)},
 \qquad m_\parallel=\gamma_0^3m.
\]
在去掉中心相位与已知线性平移的参考系中，漂移时间 $T$ 对应
\[
 T=Z_d/v_0,\qquad V(Z_d)=\ee^{-\ii \varphi_rN^2},\qquad
 \varphi_r=\frac{\hbar\omega^2Z_d}{2m_\parallel v_0^3}.
\]
这里使用固定距离的传播因子 $\ee^{\ii p_\ell Z_d/\hbar}$，二次相位来自动量对能量的非线性依赖。严格等间距 $E_\ell$ 在固定时间的传播中只产生关于 $\ell$ 的线性相位。
二阶近似还要求 $\max_{|\ell|\le M}Z_d|R_\ell^{(p)}|/\hbar\ll1$。
第一次散射后的密度矩阵因而变为
\[
 \rho^{\rm drift}_{\ell m}
 =a_\ell a_m^*\ee^{-\ii \varphi_r(\ell^2-m^2)}.
\]
取 $\ell=m$ 后相位抵消，所以漂移前后的 $P_\ell^{(1)}$ 完全相同。
非对角元仍保留相位差，所以电子的纵向相干观测量会改变，即使能谱保持不变。
\end{exbox}

\begin{exbox}[第四步：同一个态通过相位噪声通道]
设这一装置的随机相位 $\chi$ 服从均值零、方差 $2D$ 的正态分布，$D>0$。求噪声平均后的密度矩阵和相位测量概率。

\textbf{解。} 对相位噪声作高斯平均，得到
\begin{align*}
 \Phi_D(\rho)&=\int_{\mathbb R}\frac{\ee^{-\chi^2/(4D)}}{\sqrt{4\pi D}}
 \ee^{-\ii \chi N}\rho \ee^{\ii \chi N}\,\dd\chi,\\
 [\Phi_D(\rho)]_{\ell m}&=\ee^{-D(\ell-m)^2}\rho_{\ell m}.
\end{align*}
这是酉共轭的概率混合，故完全正且保迹；$D=0$ 定义为恒等通道的极限。
它描述这里指定的经典相位抖动，不自动代表任意环境的微观退相干。
结合漂移，第二次散射前的同一个态为
\[
 \rho^{\rm in}_{\ell m}
 =a_\ell a_m^*\ee^{-\ii \varphi_r(\ell^2-m^2)}\ee^{-D(\ell-m)^2}.
\]
对圆周上的协变相位测量，概率密度可以写成
\[
 p(\theta)=\frac1{2\pi}\sum_{\ell,m}
 \rho^{\rm in}_{\ell m}\ee^{-\ii (\ell-m)\theta},\qquad
 \int_0^{2\pi}p(\theta)\,\dd\theta=1.
\]
保留 $a_0=1+O(g^2),a_1=\beta+O(g^3),a_{-1}=-\beta^*+O(g^3)$，得到
\[
 p(\theta)=\frac1{2\pi}
 \left[1+4g \ee^{-D}\sin\varphi_r\sin(\phi-\theta)+O(g^2)\right].
\]
第一次脉冲刚结束时的相位密度严格均匀；漂移把相位调制转成密度调制，
而噪声压低一阶谐波。这里展开式须在 $g\ll1$ 下使用，不能把其截断形式当成任意耦合下的概率密度。
\end{exbox}

\begin{exbox}[第五步：第二次近场把相干信息转成能谱]
用第二段近场读出相干信息，求最终能谱，并检验零漂移、无噪声时的相位干涉。

\textbf{解。} 读出脉冲 $S(\beta_2)$ 的矩阵元为
\[
 U_{nm}=J_{n-m}(2g_2)\ee^{\ii (n-m)\phi_2}.
\]
所以最终能谱的完整表达是
\begin{align*}
 P_n^{\rm out}
 &=\sum_{m,m'}U_{nm}U_{nm'}^*a_ma_{m'}^*
 \ee^{-\ii \varphi_r(m^2-m'^2)}\ee^{-D(m-m')^2}.
\end{align*}
双重求和保留了不同入射能级之间的相干项，其中漂移相位与噪声因子分别由前两步确定。
若 $D=0$，它就是
\[
 P_n^{\rm out}=|A_n|^2,\qquad
 A_n=\sum_mJ_{n-m}(2g_2)\ee^{\ii (n-m)\phi_2}
 \ee^{-\ii \varphi_rm^2}J_m(2g_1)\ee^{\ii m\phi_1}.
\]
特别地，零漂移且无噪声时，两个散射生成元交换，因此
\[
 S(\beta_2)S(\beta_1)=S(\beta_1+\beta_2),\qquad
 P_n^{\rm out}=J_n^2(2|\beta_1+\beta_2|).
\]
若两脉冲等强 $g_1=g_2=g$，令 $\delta=\phi_1-\phi_2$，则
\[
 |\beta_1+\beta_2|=2g|\cos(\delta/2)|,
 \qquad P_n^{\rm out}=J_n^2\!\left(4g|\cos(\delta/2)|\right).
\]
同相时脉冲增强，反相时完全抵消。这是能量测量读出相对相位的最简单极限。
存在漂移或噪声时必须回到上面的双重求和，不能继续套用这条抵消公式。
\end{exbox}

\begin{exbox}[第六步：有限分辨率与漏检写成同一组 POVM]
将有限能量分辨率和探测效率写成测量算符，求各记录窗及漏检结果的概率。

\textbf{解。} 令仪器读出的能量变化为 $\varepsilon$，高斯响应为
\[
 R(\varepsilon|n)=\frac{1}{\sqrt{2\pi}\sigma_E}
 \exp\!\left[-\frac{(\varepsilon-n\hbar\omega)^2}{2\sigma_E^2}\right],
 \qquad F_\varepsilon=\sum_nR(\varepsilon|n)|n\rangle\langle n|.
\]
因为 $\int R(\varepsilon|n)\dd\varepsilon=1$，有 $\int F_\varepsilon\dd\varepsilon=I$。
对不交叠的记录窗 $W_j$，定义
\[
 r_{jn}=\int_{W_j}R(\varepsilon|n)\,\dd\varepsilon,\qquad
 F_j=\eta_j\sum_nr_{jn}|n\rangle\langle n|,
 \quad 0\le\eta_j\le1.
\]
于是记录概率为
\[
 p_j=\operatorname{Tr}(F_j\rho^{\rm out})
 =\eta_j\sum_nr_{jn}P_n^{\rm out}.
\]
有限窗口或效率低于一时，必须补上
\[
 F_{\rm miss}=I-\sum_jF_j\ge0,\qquad
 p_{\rm miss}=1-\sum_jp_j.
\]
这一步把能级混叠与漏检分开了。$\sigma_E\ll\hbar\omega$ 时可以近似分辨边带；
当窗口混叠明显时，真实边带占据不能直接等同于该窗口计数比例。
\end{exbox}

\begin{exbox}[第七步：计数似然、相位可辨识性与一个数值设计]
由记录概率构造计数似然，判断单个参考相位设置能否确定相对相位。

\textbf{解。} 给定固定入射电子数 $N_{\rm inc}$，把漏检也作为结果，全部计数服从多项分布。
这里要求各电子独立且同态制备，相位噪声按电子独立采样。
若一批电子共享同一个慢变化相位，单电子通道仍可相同，批计数却是条件分布的混合，不能直接沿用独立分流。
若入射数本身为均值 $\Lambda$ 的泊松变量，各窗口分流后得到独立泊松计数。再加上相互独立、已标定均值 $B_j$ 的泊松背景，有
\[
 C_j\sim\operatorname{Poisson}(\lambda_j),\qquad
 \lambda_j(\delta)=\Lambda p_j(\delta)+B_j.
\]
似然、得分和 Fisher 信息依次为
\begin{align*}
 \log\mathcal L(\delta)&=\sum_j[C_j\log\lambda_j-\lambda_j-\log(C_j!)],\\
 \partial_\delta\log\mathcal L&=\sum_j
 \left(\frac{C_j}{\lambda_j}-1\right)\partial_\delta\lambda_j,\\
 \mathcal I(\delta)&=\sum_{j:\lambda_j>0}
 \frac{(\partial_\delta\lambda_j)^2}{\lambda_j}.
\end{align*}
得分与通常的 Fisher 公式在所用窗口均有 $\lambda_j>0$ 的正则参数邻域成立。
理想反相、无背景的暗窗口满足 $\lambda_j=0$，此时应直接使用原始似然及相应极限；
不能将 $0/0$ 代入，也不能因删去该项就宣布暗端没有参数信息。
只有第一次脉冲并直接测能量时，$p_j$ 与 $\phi_1$ 无关，故这项相位信息严格为零。
读出脉冲引入后，零漂移、无噪声、等强弱脉冲的两条一阶边带合概率为
\[
 q(\delta)=P_1^{\rm out}+P_{-1}^{\rm out}
 =4g^2(1+\cos\delta)+O(g^4).
\]
一次读出仍只看到 $\cos\delta$，不能区分 $\delta$ 与 $-\delta$。
将已知参考相位增加 $\pi/2$ 后重复实验，两个设置分别测得余弦与正弦分量，
才可以在模 $2\pi$ 意义下消除这项歧义。

取 $g=0.10$、每个设置 $N_{\rm inc}=10^4$、理想分辨率、相同效率 $\eta=0.8$，忽略背景。
按弱耦合近似，$\delta=0,\pi/2,\pi$ 的两边带期望计数依次为
\[
 \eta N_{\rm inc}q(\delta)\simeq640,\quad320,\quad0.
\]
这是设计期望，并非虚构的实测数据。精确预测应使用 Bessel 概率；
其中反相设置在当前理想极限下严格为零，另两个数含弱耦合截断误差。
\end{exbox}

\begin{exbox}[第八步：反演场幅与控制计算窗]
由第一次散射的计数估计场幅，并用边带尾部估计有限计算窗的截断误差。

\textbf{解。} 先关闭读出脉冲，对第一次散射的两条一阶边带做理想分辨率标定。
在弱耦合、已知效率 $\eta$、无背景条件下，边带总计数 $C$ 近似满足
\[
 C\sim\operatorname{Binomial}(N_{\rm inc},2\eta g^2),\qquad
 \widehat g=\sqrt{\frac{C}{2\eta N_{\rm inc}}}.
\]
它是弱耦合模型的反演估计，并非精确 Bessel 模型下的无偏估计。
当期望计数足够大，Delta 方法给出
\[
 \operatorname{Var}(\widehat g)
 \simeq\frac{1-2\eta g^2}{8\eta N_{\rm inc}},\qquad
 |\widehat E_a|=\frac{\hbar\omega\widehat g}{|q_e|\sqrt{2\pi}L}
 \ee^{L^2(\kappa-\omega/v_0)^2/2}.
\]
因此大的相位失配既压低信号，也放大场幅反演的不确定性。
若有背景、窗口混叠或未知效率，就必须使用第六、七步的完整概率模型联合拟合；
不能只减去平均背景后仍沿用无背景二项方差。

对第一次散射，第二步算出的二阶矩还给出可直接用于数值截断的界
\[
 \epsilon_M=\sum_{|\ell|>M}|a_\ell|^2
 \le\frac{2g^2}{(M+1)^2}.
\]
归一化截断纯态与原态的迹距离为 $\sqrt{\epsilon_M}$，
故在后续使用相同的完全正保迹演化和相同测量时，结果概率的总变差不超过这一界。
这不包括把第二次脉冲自身替换为有限矩阵所增加的边界误差；
也不自动控制能量这类无界观测量的期望，后者需要另行估计尾部矩。

\end{exbox}

场剖面的 $E_a,L,\kappa,v_0$ 先决定 $\beta_1$，随后得到的电子振幅经历漂移、噪声和参考脉冲，最终转成记录概率与计数似然。因而场幅反演依赖整条演化和测量链。色散截断、相位噪声模型、仪器响应和计数采样各有自己的使用条件；改变其中一项，就应从对应的那一步重新计算预测与不确定性。
